\documentclass[tikz,border=10pt]{standalone}
\usepackage{tikz}
\usepackage{amsmath}
\usepackage{amssymb}
\usepackage{ctex}
\usepackage{pgfplots}
\pgfplotsset{compat=1.18}
\usetikzlibrary{calc, positioning, arrows.meta, decorations.pathreplacing}

\begin{document}
\begin{tikzpicture}[>=Stealth]

%% ============================
%%  左子图：正定 — f=x^2+y^2，椭圆等高线（碗状）
%%  特征值 λ1=λ2=1 > 0
%% ============================
\begin{scope}[xshift=0cm]
  \begin{axis}[
    name=ax1,
    width=6.2cm, height=6.2cm,
    xmin=-2.2, xmax=2.2,
    ymin=-2.2, ymax=2.2,
    axis lines=center,
    xlabel={$x$}, ylabel={$y$},
    xlabel style={font=\small, right},
    ylabel style={font=\small, above},
    xtick={-2,-1,0,1,2}, ytick={-2,-1,0,1,2},
    tick label style={font=\scriptsize},
    grid=both, grid style={gray!12},
    clip=true, axis equal image,
  ]
    % 等高线: x^2+y^2 = c，圆族
    \addplot[blue!80!black, thick, domain=0:360, samples=100]
      ({0.5*cos(x)}, {0.5*sin(x)});
    \addplot[blue!70!black, thick, domain=0:360, samples=100]
      ({1.0*cos(x)}, {1.0*sin(x)});
    \addplot[blue!60!black, thick, domain=0:360, samples=100]
      ({1.5*cos(x)}, {1.5*sin(x)});
    \addplot[blue!50!black, thick, domain=0:360, samples=100]
      ({2.0*cos(x)}, {2.0*sin(x)});

    % 等高线标注
    \node[font=\scriptsize,blue!80!black] at (axis cs:0.38,0.2) {$c_1$};
    \node[font=\scriptsize,blue!70!black] at (axis cs:0.78,0.55) {$c_2$};
    \node[font=\scriptsize,blue!60!black] at (axis cs:1.18,0.82) {$c_3$};

    % 最小值中心
    \addplot[only marks, mark=*, mark size=3pt, red!70!black]
      coordinates {(0,0)};
    \node[font=\scriptsize, red!70!black, above right]
      at (axis cs:0.05,0.05) {最小值};

    % 特征值
    \node[font=\scriptsize, blue!80!black,
          fill=blue!10, draw=blue!40, rounded corners=2pt, inner sep=3pt,
          align=center]
      at (axis cs:0, 1.9) {
        $\lambda_1=1>0$\\
        $\lambda_2=1>0$
      };
  \end{axis}

  \node[font=\small\bfseries, blue!80!black,
        fill=blue!10, draw=blue!50, rounded corners=4pt,
        inner sep=5pt, align=center]
    at ([yshift=0.4cm]ax1.north) {\textbf{正定}（Positive Definite）};

  \node[font=\small, align=center] at ([yshift=-0.8cm]ax1.south)
    {$f = x^2+y^2$（碗状）\\
     等高线 $=$ 圆族，全局极小值};
\end{scope}

%% ============================
%%  中子图：鞍点 — f=x^2-y^2，双曲等高线（X字形）
%%  特征值 λ1=1>0, λ2=-1<0
%% ============================
\begin{scope}[xshift=8.0cm]
  \begin{axis}[
    name=ax2,
    width=6.2cm, height=6.2cm,
    xmin=-2.2, xmax=2.2,
    ymin=-2.2, ymax=2.2,
    axis lines=center,
    xlabel={$x$}, ylabel={$y$},
    xlabel style={font=\small, right},
    ylabel style={font=\small, above},
    xtick={-2,-1,0,1,2}, ytick={-2,-1,0,1,2},
    tick label style={font=\scriptsize},
    grid=both, grid style={gray!12},
    clip=true, axis equal image,
  ]
    % 等高线 x^2-y^2=c (双曲线)
    % c>0 右支+左支: x1=sqrt(c+y^2) => 参数化: x=sqrt(c)*cosh(t), y=sqrt(c)*sinh(t)
    % 避免 cosh/sinh 除零，用参数 t 直接数值
    % c=1: x^2-y^2=1 右支 x=sqrt(1+y^2)
    \addplot[purple!70!black, thick, domain=-2.0:2.0, samples=80]
      ({sqrt(1.0 + x*x)}, {x});
    \addplot[purple!70!black, thick, domain=-2.0:2.0, samples=80]
      ({-sqrt(1.0 + x*x)}, {x});

    % c=-1: x^2-y^2=-1 => y^2-x^2=1 上支/下支
    \addplot[purple!50!black, thick, domain=-2.0:2.0, samples=80]
      ({x}, {sqrt(1.0 + x*x)});
    \addplot[purple!50!black, thick, domain=-2.0:2.0, samples=80]
      ({x}, {-sqrt(1.0 + x*x)});

    % c=0.25 右支
    \addplot[purple!80!black, thick, domain=-2.0:2.0, samples=80]
      ({sqrt(0.25 + x*x)}, {x});
    \addplot[purple!80!black, thick, domain=-2.0:2.0, samples=80]
      ({-sqrt(0.25 + x*x)}, {x});

    % 渐近线 y=±x （X字形）
    \addplot[gray!50, dashed, thick, domain=-2.0:2.0] {x};
    \addplot[gray!50, dashed, thick, domain=-2.0:2.0] {-x};

    % 鞍点
    \addplot[only marks, mark=*, mark size=3pt, orange!80!black]
      coordinates {(0,0)};
    \node[font=\scriptsize, orange!80!black, above right]
      at (axis cs:0.08,0.08) {鞍点};

    % 特征值
    \node[font=\scriptsize, purple!70!black,
          fill=purple!8, draw=purple!40, rounded corners=2pt, inner sep=3pt,
          align=center]
      at (axis cs:0, -1.7) {
        $\lambda_1=1>0$\\
        $\lambda_2=-1<0$
      };
  \end{axis}

  \node[font=\small\bfseries, purple!70!black,
        fill=purple!10, draw=purple!50, rounded corners=4pt,
        inner sep=5pt, align=center]
    at ([yshift=0.4cm]ax2.north) {\textbf{不定（鞍点）}};

  \node[font=\small, align=center] at ([yshift=-0.8cm]ax2.south)
    {$f = x^2-y^2$（鞍面）\\
     等高线 $=$ 双曲线 X 字形};
\end{scope}

%% ============================
%%  右子图：半正定 — f=x^2，只对x二次（槽状）
%%  特征值 λ1=1>0, λ2=0
%% ============================
\begin{scope}[xshift=16.0cm]
  \begin{axis}[
    name=ax3,
    width=6.2cm, height=6.2cm,
    xmin=-2.2, xmax=2.2,
    ymin=-2.2, ymax=2.2,
    axis lines=center,
    xlabel={$x$}, ylabel={$y$},
    xlabel style={font=\small, right},
    ylabel style={font=\small, above},
    xtick={-2,-1,0,1,2}, ytick={-2,-1,0,1,2},
    tick label style={font=\scriptsize},
    grid=both, grid style={gray!12},
    clip=true, axis equal image,
  ]
    % 等高线 x^2=c → 竖线 x=±sqrt(c)
    \addplot[teal!80!black, thick] coordinates {(0.5,-2.2) (0.5,2.2)};
    \addplot[teal!80!black, thick] coordinates {(-0.5,-2.2) (-0.5,2.2)};

    \addplot[teal!70!black, thick] coordinates {(1.0,-2.2) (1.0,2.2)};
    \addplot[teal!70!black, thick] coordinates {(-1.0,-2.2) (-1.0,2.2)};

    \addplot[teal!60!black, thick] coordinates {(1.5,-2.2) (1.5,2.2)};
    \addplot[teal!60!black, thick] coordinates {(-1.5,-2.2) (-1.5,2.2)};

    % 零等高线 x=0（y轴）
    \addplot[teal!90!black, very thick] coordinates {(0,-2.2) (0,2.2)};
    \node[font=\scriptsize, teal!80!black, right]
      at (axis cs:0.05, 1.5) {$f=0$};

    % 标注方向
    \draw[->, teal!60, thick] (axis cs:0,0) -- (axis cs:1.5,0)
      node[below, font=\scriptsize] {$f$ 增大};

    % 特征值
    \node[font=\scriptsize, teal!80!black,
          fill=teal!8, draw=teal!40, rounded corners=2pt, inner sep=3pt,
          align=center]
      at (axis cs:0, -1.7) {
        $\lambda_1=1>0$\\
        $\lambda_2=0$（退化）
      };
  \end{axis}

  \node[font=\small\bfseries, teal!80!black,
        fill=teal!10, draw=teal!50, rounded corners=4pt,
        inner sep=5pt, align=center]
    at ([yshift=0.4cm]ax3.north) {\textbf{半正定}（Positive Semi-definite）};

  \node[font=\small, align=center] at ([yshift=-0.8cm]ax3.south)
    {$f = x^2$（槽状）\\
     等高线 $=$ 平行竖线（$y$ 方向退化）};
\end{scope}

%% ===== 底部判定准则框 =====
\node[font=\small, fill=white, draw=gray!50, thin, rounded corners=4pt,
      inner sep=10pt, align=center]
  at (11.2, -3.2) {
    \textbf{二次型 $\mathbf{x}^TA\mathbf{x}$ 等高线形状 $\Leftrightarrow$ 特征值符号：}
    \quad 正定（所有 $\lambda>0$）$\to$ 椭圆
    \quad 不定（$\lambda$ 有正有负）$\to$ 双曲线
    \quad 半正定（$\lambda\geq0$，含零）$\to$ 平行线
  };

%% ===== 整体标题 =====
\node[font=\large\bfseries] at (11.2, 4.2)
  {二次型曲面俯视等高线};

\end{tikzpicture}
\end{document}
